\chapter{Analyse et conception}
\label{chap:conception}

Ce chapitre formalise la conception de \devpilot{} : identification des acteurs,
cas d'utilisation, modèle du domaine, diagrammes de classes, de séquence, modèle
de données, et vues d'architecture (logique, physique, déploiement, composants,
paquets, communication). Une attention particulière est portée à la modélisation
des quatre modules distinctifs : le moteur RAG, RAGOps, le système de traçabilité
et l'AI Execution Studio. Chaque diagramme est introduit par son explication, ses
hypothèses de modélisation, son code PlantUML et son interprétation.

\section{Identification des acteurs}
Trois acteurs humains interagissent avec le système, selon une hiérarchie de
droits : l'\textbf{Utilisateur} (membre d'une équipe d'ingénierie), l'\textbf{Administrateur}
(gestion de la plateforme) et le \textbf{Super-administrateur} (observabilité
approfondie). Trois acteurs systèmes externes participent au pipeline :
\textbf{Qdrant} (base vectorielle), \textbf{Ollama} (\emph{embeddings}) et
\textbf{Gemini} (génération).

\section{Cas d'utilisation}
Les cas d'utilisation clés sont : gérer un espace de travail, téléverser et suivre
un document, poser une question (RAG) et recevoir une réponse citée en
\emph{streaming}, consulter l'historique, superviser les pipelines (RAGOps) et
inspecter une trace. Le tableau~\ref{tab:uc} détaille le cas pivot.

\begin{table}[H]
\centering\caption{Cas d'utilisation détaillé : « Poser une question (RAG) »}
\label{tab:uc}\small
\begin{tabularx}{\textwidth}{@{}>{\bfseries}l X@{}}
\toprule
Acteur & Utilisateur authentifié, membre d'un espace de travail \\
Pré-condition & L'espace contient au moins un document indexé \\
Déclencheur & Soumission d'une question dans l'interface de chat \\
Scénario nominal & Router, réécriture, récupération, reclassement, génération, évaluation, correction, diffusion en \emph{streaming} \\
Post-condition & Message, agents, passages et évaluation persistés \\
Exceptions & Aucun passage pertinent $\rightarrow$ repli explicite ; échec LLM $\rightarrow$ modèle de \emph{fallback} \\
\bottomrule
\end{tabularx}
\end{table}

\section{Diagramme de cas d'utilisation}
\paragraph{Explication.} Ce diagramme présente les acteurs et leurs interactions
avec les fonctionnalités principales, ainsi que la relation d'héritage entre rôles.
\paragraph{Hypothèses.} Les acteurs systèmes (Qdrant, Ollama, Gemini) sont
représentés comme acteurs secondaires sollicités par le cas « Poser une question ».
Le super-administrateur hérite de l'administrateur, qui hérite de l'utilisateur.

\begin{lstlisting}[style=plantuml,caption={Diagramme de cas d'utilisation},label={lst:uc}]
@startuml
left to right direction
actor Utilisateur as U
actor Administrateur as A
actor "Super-administrateur" as S
actor Qdrant as Q
actor Gemini as G

A --|> U
S --|> A

rectangle "DevPilot AI" {
  usecase "Gérer un espace" as UC1
  usecase "Téléverser un document" as UC2
  usecase "Poser une question (RAG)" as UC3
  usecase "Consulter l'historique" as UC4
  usecase "Administrer les utilisateurs" as UC5
  usecase "Superviser (RAGOps)" as UC6
  usecase "Inspecter une trace" as UC7
}

U --> UC1
U --> UC2
U --> UC3
U --> UC4
A --> UC5
A --> UC6
S --> UC7
UC3 ..> Q : <<include>>
UC3 ..> G : <<include>>
@enduml
\end{lstlisting}

\paragraph{Interprétation.} L'utilisateur dispose des cas métier (espaces,
documents, RAG, historique). L'administrateur ajoute la supervision, le
super-administrateur l'inspection fine. Le cas RAG inclut systématiquement la
recherche vectorielle et la génération.
\plantumlfig{Diagramme de cas d'utilisation global}{fig:uc}

\section{Modèle du domaine}
\paragraph{Explication.} Le modèle du domaine articule les concepts métier
centraux et leurs relations, indépendamment de la persistance.
\paragraph{Hypothèses.} La \emph{Trace} est élevée au rang de concept du domaine
(agrégat des exécutions d'agents, passages récupérés et évaluation), choix
déterminant pour l'observabilité.

\begin{lstlisting}[style=plantuml,caption={Modèle du domaine},label={lst:domaine}]
@startuml
class Utilisateur
class EspaceTravail
class Adhesion
class Document
class Chunk
class Conversation
class Message
class Trace

Utilisateur "1" -- "0..*" Adhesion
EspaceTravail "1" -- "0..*" Adhesion
EspaceTravail "1" -- "0..*" Document
Document "1" *-- "0..*" Chunk
Utilisateur "1" -- "0..*" Conversation
Conversation "1" *-- "1..*" Message
Message "1" -- "0..1" Trace
@enduml
\end{lstlisting}

\paragraph{Interprétation.} Un utilisateur accède aux espaces via des adhésions ;
un espace contient des documents découpés en \emph{chunks} ; les conversations
agrègent des messages, et chaque message assistant porte une trace.
\plantumlfig{Modèle du domaine}{fig:domaine}

\section{Diagrammes de classes}
\subsection{Entités persistantes}
\paragraph{Explication.} Ce diagramme présente les onze entités persistantes avec
leurs attributs clés et leurs cardinalités.
\paragraph{Hypothèses.} Les identifiants sont des UUID ; les attributs techniques
communs (horodatages) sont omis pour la lisibilité.

\begin{lstlisting}[style=plantuml,caption={Diagramme de classes — entités persistantes},label={lst:classes}]
@startuml
class users { +id : UUID\n+email\n+role\n+is_active }
class workspaces { +id : UUID\n+name\n+owner_id }
class workspace_members { +id : UUID\n+role }
class documents { +id : UUID\n+filename\n+status }
class chunks { +id : UUID\n+chunk_index\n+vector_id }
class conversations { +id : UUID\n+title }
class messages { +id : UUID\n+role\n+content }
class agent_runs { +id : UUID\n+agent_type\n+latency_ms }
class retrieved_chunks { +id : UUID\n+rank\n+score }
class evaluations { +id : UUID\n+faithfulness\n+hallucination_score }
class llm_usage { +id : UUID\n+total_tokens\n+cost_usd }

users "1" -- "0..*" workspaces
users "1" -- "0..*" workspace_members
workspaces "1" -- "0..*" workspace_members
workspaces "1" *-- "0..*" documents
documents "1" *-- "0..*" chunks
users "1" -- "0..*" conversations
conversations "1" *-- "1..*" messages
messages "1" -- "0..*" agent_runs
messages "1" -- "0..*" retrieved_chunks
messages "1" -- "1" evaluations
messages "1" -- "0..*" llm_usage
@enduml
\end{lstlisting}

\paragraph{Interprétation.} Les compositions (\emph{chunks} dans \emph{documents},
\emph{messages} dans \emph{conversations}) traduisent un cycle de vie lié ; la
relation 1--1 entre \texttt{messages} et \texttt{evaluations} formalise l'unicité
de l'évaluation par réponse.
\plantumlfig{Diagramme de classes des entités persistantes}{fig:classes}

\subsection{Moteur agentique}
\paragraph{Explication.} Ce diagramme modélise le flux agentique : un orchestrateur
coordonne sept agents, appuyés par des services applicatifs.
\paragraph{Hypothèses.} Les agents partagent une interface commune d'exécution ;
seuls les services principaux sont représentés.

\begin{lstlisting}[style=plantuml,caption={Diagramme de classes — moteur agentique},label={lst:agents}]
@startuml
class AgenticRAGWorkflow {
  +run(question) : Result
}
interface Agent {
  +execute(ctx)
}
class RouterAgent
class QueryRewriterAgent
class RetrievalAgent
class RerankerAgent
class GeneratorAgent
class EvaluatorAgent
class CorrectorAgent
class RetrievalService
class RerankingService
class EvaluationService
class RagTracesService

RouterAgent ..|> Agent
QueryRewriterAgent ..|> Agent
RetrievalAgent ..|> Agent
RerankerAgent ..|> Agent
GeneratorAgent ..|> Agent
EvaluatorAgent ..|> Agent
CorrectorAgent ..|> Agent

AgenticRAGWorkflow o-- RouterAgent
AgenticRAGWorkflow o-- QueryRewriterAgent
AgenticRAGWorkflow o-- RetrievalAgent
AgenticRAGWorkflow o-- RerankerAgent
AgenticRAGWorkflow o-- GeneratorAgent
AgenticRAGWorkflow o-- EvaluatorAgent
AgenticRAGWorkflow o-- CorrectorAgent
RetrievalAgent --> RetrievalService
RerankerAgent --> RerankingService
EvaluatorAgent --> EvaluationService
AgenticRAGWorkflow --> RagTracesService
@enduml
\end{lstlisting}

\paragraph{Interprétation.} L'orchestrateur agrège les sept agents (agrégation
\texttt{o--}) qui réalisent une interface commune ; les agents délèguent aux
services métier la récupération, le reclassement et l'évaluation.
\plantumlfig{Diagramme de classes du moteur agentique}{fig:agents}

\section{Diagrammes de séquence}
\subsection{Ingestion documentaire asynchrone}
\paragraph{Explication.} Séquence du téléversement à l'indexation, exécutée hors
du chemin critique par Celery.
\paragraph{Hypothèses.} Le stockage de fichiers et la base relationnelle sont
distingués ; la réponse HTTP est immédiate (202).

\begin{lstlisting}[style=plantuml,caption={Séquence — ingestion documentaire},label={lst:seq-ingest}]
@startuml
actor Client
participant "API (FastAPI)" as API
participant "Stockage" as ST
participant "Celery" as CEL
participant "Qdrant" as QD

Client -> API : POST /documents (fichier)
activate API
API -> ST : persister fichier + métadonnées
API -> CEL : enqueue process_document_task
API --> Client : 202 Accepted (statut=pending)
deactivate API
activate CEL
CEL -> ST : lire + extraire le texte
CEL -> CEL : découper + vectoriser (768-d)
CEL -> QD : upsert vecteurs
CEL -> ST : statut = indexed
deactivate CEL
@enduml
\end{lstlisting}

\paragraph{Interprétation.} L'API rend la main immédiatement ; le \emph{worker}
réalise l'extraction, le découpage, la vectorisation et l'indexation, puis met à
jour le statut du document.
\plantumlfig{Diagramme de séquence — ingestion asynchrone}{fig:seq-ingest}

\subsection{Requête RAG en streaming}
\paragraph{Explication.} Séquence de la requête RAG diffusée en \emph{streaming}
SSE à travers le flux agentique.
\paragraph{Hypothèses.} Les sept agents sont synthétisés au niveau du
\texttt{Workflow} ; les événements SSE sont émis au fil de l'eau.

\begin{lstlisting}[style=plantuml,caption={Séquence — requête RAG en streaming},label={lst:seq-rag}]
@startuml
actor Client
participant "API (SSE)" as API
participant "Workflow" as WF
participant "Qdrant" as QD
participant "Gemini" as LLM

Client -> API : POST /chat/stream (question)
activate API
API -> WF : exécuter le pipeline
activate WF
WF -> WF : router + réécriture
WF -> QD : récupérer candidats (k=15)
QD --> WF : passages + scores
WF -> WF : reclasser (top 5)
API --> Client : event: trace / citations
WF -> LLM : générer (stream)
LLM --> WF : flux de jetons
API --> Client : event: token (xN)
WF -> WF : évaluer + corriger
API --> Client : event: evaluation / message / done
deactivate WF
deactivate API
@enduml
\end{lstlisting}

\paragraph{Interprétation.} La diffusion entrelace les étapes du pipeline et les
événements SSE : les citations précèdent les jetons, suivis de l'évaluation et de
l'événement de clôture.
\plantumlfig{Diagramme de séquence — requête RAG en streaming}{fig:seq-rag}

\subsection{Inspection d'une trace}
\paragraph{Explication.} Séquence de consultation d'une trace par
l'administrateur.
\paragraph{Hypothèses.} L'inspecteur agrège quatre appels parallèles (trace,
récupération, reclassement, évaluation).

\begin{lstlisting}[style=plantuml,caption={Séquence — inspection de trace},label={lst:seq-trace}]
@startuml
actor Admin
participant "Frontend" as FE
participant "API Admin" as API
database "PostgreSQL" as PG

Admin -> FE : ouvrir /admin/rag-traces/{id}
FE -> API : GET trace + retrieval + reranking + evaluation
activate API
API -> PG : lire la lignée du message
PG --> API : agent_runs, retrieved_chunks, evaluations
API --> FE : trace complète
deactivate API
FE --> Admin : reconstruction des 9 sections
@enduml
\end{lstlisting}

\paragraph{Interprétation.} La trace est reconstruite intégralement à partir des
données persistées, sans dépendance à des journaux externes.
\plantumlfig{Diagramme de séquence — inspection de trace}{fig:seq-trace}

\section{Modèle de données}
\paragraph{Explication.} Le modèle relationnel comprend onze tables ; la
figure~\ref{fig:erd} en présente la vue entité-association.
\paragraph{Hypothèses.} Les clés étrangères sont implicites dans les associations ;
les colonnes techniques sont omises.

\begin{lstlisting}[style=plantuml,caption={Modèle de données (ERD)},label={lst:erd}]
@startuml
entity users
entity workspaces
entity workspace_members
entity documents
entity chunks
entity conversations
entity messages
entity agent_runs
entity retrieved_chunks
entity evaluations
entity llm_usage

users ||--o{ workspaces
users ||--o{ workspace_members
workspaces ||--o{ workspace_members
workspaces ||--o{ documents
documents ||--o{ chunks
users ||--o{ conversations
conversations ||--o{ messages
messages ||--o{ agent_runs
messages ||--o{ retrieved_chunks
messages ||--|| evaluations
messages ||--o{ llm_usage
@enduml
\end{lstlisting}

\paragraph{Interprétation.} Le schéma matérialise la chaîne de traçabilité
\texttt{message} $\rightarrow$ \texttt{agent\_runs} $\rightarrow$
\texttt{retrieved\_chunks} $\rightarrow$ \texttt{evaluations}, complétée par
\texttt{llm\_usage}.
\plantumlfig{Modèle de données — vue entité-association}{fig:erd}

\section{Architecture logique}
L'application s'organise en couches aux responsabilités strictes : présentation
(\emph{frontend}), application (API et services), domaine IA (pipeline agentique),
traitement asynchrone (Celery) et données (PostgreSQL, Qdrant, stockage). Les
dépendances ne circulent que du haut vers le bas, condition de la maintenabilité.

\section{Architecture physique}
Les composants sont déployés sous forme de conteneurs indépendants communiquant
par le réseau : services applicatifs sans état, magasins de données persistants et
services d'IA (Ollama local, Gemini distant). Cette séparation autorise une montée
en charge ciblée de chaque composant.

\section{Architecture de déploiement}
\paragraph{Explication.} Diagramme de déploiement des huit services conteneurisés,
de l'hôte Ollama et de l'API Gemini externe.
\paragraph{Hypothèses.} Un n{\oe}ud unique héberge les conteneurs ; Ollama
s'exécute sur l'hôte, Gemini est un service \emph{cloud}.

\begin{lstlisting}[style=eraser,caption={Diagramme de déploiement (Eraser)},label={lst:deploy}]
colorMode bold
styleMode shadow
direction down

Docker Host [icon: docker, color: blue] {
  frontend [icon: nextjs, label: "frontend (3001)"]
  backend [icon: python, label: "backend uvicorn (8000)"]
  celery_worker [icon: python]
  postgres [icon: postgresql, label: "postgres:16"]
  redis [icon: redis, label: "redis:7"]
  qdrant [icon: database, label: "qdrant 1.12"]
  prometheus [icon: prometheus]
  grafana [icon: grafana]
}
Ollama Host [icon: server, color: teal] { ollama [icon: cpu, label: "Ollama :11434"] }
Gemini Cloud [icon: google, color: orange] { gemini [icon: cloud, label: "Gemini API"] }

frontend > backend: HTTPS / SSE
backend > postgres
backend > redis
backend > qdrant
backend > ollama
backend > gemini
celery_worker > redis
celery_worker > postgres
celery_worker > qdrant
backend > prometheus
prometheus > grafana
\end{lstlisting}

\paragraph{Interprétation.} Le \emph{backend} et le \emph{worker} partagent les
magasins de données ; Redis relie l'API au \emph{worker} ; l'observabilité repose
sur Prometheus et Grafana.
\eraserfig{Diagramme de déploiement}{fig:deploy}

\section{Diagramme de composants}
\paragraph{Explication.} Vue des composants applicatifs et de leurs dépendances.
\paragraph{Hypothèses.} Les interfaces réseau (REST, SSE) sont représentées par
les flèches de dépendance.

\begin{lstlisting}[style=eraser,caption={Diagramme de composants (Eraser)},label={lst:comp}]
colorMode bold
styleMode shadow
direction down

Frontend [icon: nextjs, color: blue]
API [icon: python, color: blue, label: "API FastAPI"]
Workflow [icon: cpu, color: purple, label: "Pipeline agentique"]
Celery [icon: refresh-cw, color: orange]
Redis [icon: redis, color: orange]
Postgres [icon: postgresql, color: green]
Qdrant [icon: database, color: green]
Ollama [icon: server, color: teal]
Gemini [icon: google, color: teal]
Observability [icon: activity, color: red, label: "Prometheus / Grafana"]

Frontend > API
API > Workflow
API > Postgres
API > Redis
Redis > Celery
Workflow > Qdrant
Workflow > Ollama
Workflow > Gemini
Celery > Postgres
Celery > Qdrant
API > Observability
\end{lstlisting}

\paragraph{Interprétation.} Le \emph{frontend} ne dépend que de l'API ; le
pipeline agentique concentre les dépendances vers Qdrant et les fournisseurs d'IA.
\eraserfig{Diagramme de composants}{fig:comp}

\section{Diagramme de paquets}
\paragraph{Explication.} Organisation modulaire du code \emph{backend} et
\emph{frontend}.
\paragraph{Hypothèses.} Seuls les paquets principaux sont représentés.

\begin{lstlisting}[style=plantuml,caption={Diagramme de paquets},label={lst:pkg}]
@startuml
package "backend.app" {
  package api
  package core
  package db
  package models
  package schemas
  package services
  package agents
  package workers
  package utils
}
package "frontend" {
  package app
  package components
  package features
  package lib
}
api --> services
services --> models
services --> agents
agents --> utils
workers --> services
app --> features
features --> components
features --> lib
@enduml
\end{lstlisting}

\paragraph{Interprétation.} Le découpage en paquets reflète la séparation des
responsabilités : routes, logique métier, agents et accès aux données côté
serveur ; pages, modules fonctionnels, composants et bibliothèques côté client.
\plantumlfig{Diagramme de paquets}{fig:pkg}

\section{Diagramme de communication}
\paragraph{Explication.} Vue de collaboration numérotée de la requête RAG.
\paragraph{Hypothèses.} Les messages sont ordonnés selon la séquence
d'exécution.

\begin{lstlisting}[style=plantuml,caption={Diagramme de communication — requête RAG},label={lst:comm}]
@startuml
object Client
object API
object Workflow
object Qdrant
object Gemini

Client -> API : 1: question
API -> Workflow : 2: exécuter
Workflow -> Qdrant : 3: récupérer
Qdrant -> Workflow : 4: passages
Workflow -> Gemini : 5: générer
Gemini -> Workflow : 6: jetons
Workflow -> API : 7: trace + réponse
API -> Client : 8: stream + citations
@enduml
\end{lstlisting}

\paragraph{Interprétation.} La numérotation explicite l'ordre des échanges et met
en évidence le rôle central de l'orchestrateur dans la coordination.
\plantumlfig{Diagramme de communication}{fig:comm}

\section{Conclusion}
La conception traduit les besoins en modèles exploitables : cas d'utilisation,
modèle du domaine, diagrammes de classes (entités et moteur agentique), scénarios
d'interaction, modèle de données à onze tables et vues d'architecture. La
modélisation explicite de la trace et du flux agentique prépare directement
l'architecture détaillée au chapitre~\ref{chap:archi}.
